\documentclass[a4paper,oneside,12pt]{amsart}

\usepackage[spanish, activeacute]{babel}
\usepackage[utf8]{inputenc}
%\usepackage[latin1]{inputenc}
\usepackage{latexsym}
\usepackage{amssymb}
\usepackage{multicol}
\usepackage{enumerate}
\usepackage{booktabs}
\usepackage[heightrounded]{geometry}
\geometry{top=2cm,bottom=2cm,left=2cm,right=2cm}
\usepackage{graphicx}

\DeclareMathOperator{\cond}{cond}
\DeclareMathOperator{\Id}{Id}
\DeclareMathOperator{\re}{Re}
%\DeclareMathOperator{\im}{Im}
\newcommand{\I}{\mathbb{I}}
\newcommand{\R}{\mathbb{R}}
\newcommand{\Q}{\mathbb{Q}}
\newcommand{\C}{\mathbb{C}}
\newcommand{\Z}{\mathbb{Z}}
\newcommand{\N}{\mathbb{N}}
\DeclareMathOperator{\im}{Im}

\newcommand{\nombremateria}
{Investigaci\'on Operativa}
\newcommand{\nombrecuatrimestre}
    {Segundo cuatrimestre de 2018}
\newcommand{\numeroparcial}
    { Primer Parcial }
\newcommand{\fecha}
    {18 de Septiembre de 2020}
\newcommand{\numerotema}
    {\bf 1}
\newcommand{\dsum}{\displaystyle\sum}
\newcommand{\mc}[1]{\mathcal{#1}}
\newcommand{\fa}[2]{\forall\,{#1}\in #2\;}

%\oddsidemargin -1cm \evensidemargin -1cm \textwidth 17.5cm
%topmargin 1.2cm \textheight 25cm
%\setlength{\textheight}{25cm}

\def \ds{\displaystyle}
\newtheorem{quest}{Ejercicio}
\thispagestyle{empty}

\begin{document}

\hfill \hfill
\begin{tabular}[c]{|c|c|c|c|c|}
\hline 
1 & 2 & 3 &  Calificaci\'on \\ \hline \hline
\hspace{1.2 cm} & \hspace{1.2 cm} & \hspace{1.2 cm} & \hspace{1.5 cm} \\
\hspace{1.2 cm} & \hspace{1.2 cm} & \hspace{1.2 cm} & \hspace{1.5 cm} \\ \hline

\end{tabular}

\vspace{-1cm}

\noindent Nombre y apellido:

\vspace{0.2cm}

\noindent Número de libreta:

\noindent \hrulefill 

\smallskip
\renewcommand{\baselinestretch}{1.1}

\begin{center}
	\large \textbf{\nombremateria} 
	
	\numeroparcial -- \fecha
\end{center}

\vspace{.6cm}
\renewcommand{\theenumi}{\textbf{\arabic{enumi}}}

%%% EMPIEZA EL EJERCICIO 1 %%%

\begin{quest}
    [Adaptado de \textit{Linear Programming Foundations and Extensions}, de J. Vanderbei] Una pequeña aerolínea ofrece vuelos entre tres ciudades: Córdoba, Río Gallegos y Neuquén. En este problema nos vamos a enfocar en el vuelo que se llevan a cabo los días Viernes: parte de Córdoba, hace una escala en Neuquén y luego continúa a Río Gallegos. Hay tres tipos de pasajeros:
    \begin{itemize}
        \item los que viajan de Córdoba a Neuquén.
        \item los que viajan de Neuquén a Río Gallegos.
        \item los que viajan de Córdoba a Río Gallegos.
    \end{itemize}
    El avión es pequeño y puede llevar hasta N pasajeros. La aerolínea ofrece tres tarifas distintas:
    \begin{itemize}
        \item A: dos equipajes en bodega y un equipaje de mano.
        \item B: un equipaje en bodega y un equipaje de mano.
        \item C: sólo equipaje de mano.
    \end{itemize}
    Los precios de cada pasaje ya fueron determinados y publicitados:
    \begin{table}[h!]
        \centering
        \begin{tabular}{cccc}
            Tarifa & Córdoba-Neuquén & Neuquén-Río Gallegos & Córdoba-Río Gallegos \\ \midrule
            A & 300 & 160 & 360 \\
            B & 220 & 130 & 280 \\
            C & 100 & 80 & 140 \\
        \end{tabular}
    \end{table}
    
    Por otro lado, luego de un estudio de vuelos anteriores, se han determinado las siguientes cotas inferiores en la cantidad de potenciales clientes para cada una de las combinaciones origen-destino/tarifa:
    \begin{table}[h!]
        \centering
        \begin{tabular}{cccc}
            Tarifa & Córdoba-Neuquén & Neuquén-Río Gallegos & Córdoba-Río Gallegos \\ \midrule
            A & 12 & 18 & 21 \\
            B & 18 & 28 & 30 \\
            C & 32 & 45 & 40 \\
        \end{tabular}
    \end{table}
    
    Se desea decidir cuántos pasajes vender para cada una de las nueve combinaciones origen-destino/tarifa. Tener en cuenta las siguientes restricciones:
    \begin{enumerate}
        \item la cantidad de pasajes puestos en venta para cada combinación origen-destino/tarifa no debe ser menor que la cota inferior de demanda.
        \item para ninguno de los dos tramos del vuelo pueden haber más reservas que la capacidad del avión.
        \item debido a la capacidad de carga del avión, los pasajeros de tarifa $A$ en cada tramo del viaje deben ser a lo sumo el $\%20$ de los pasajeros totales.
        \item si para Córdoba-Neuquén se pone a la venta menos de $25$ pasajes de tarifa $B$, entonces poner a la venta al menos $40$ pasajes de tarifa $C$ para Córdoba-Neuquén.
    \end{enumerate}
    Diseñar un modelo de programación lineal entera cuyo objetivo sea maximizar la ganancia.
\end{quest}

%%% TERMINA EL EJERCICIO 1 %%%

%%% EMPIEZA EL EJERCICIO 2 %%%
\begin{quest}

    Se desea organizar un festival de música que se celebrará a lo largo de tres días. Cada jornada tendrá una duración de $14$ horas divididas en franjas de dos horas, acumulando un total de $21$ franjas horarias a lo largo de los tres días. A cada banda invitada se le asignará exactamente una de esas franjas horarias para realizar su show.
    
    Se cuenta con cuatro escenarios $E_1,E_2,E_3,E_4$, por lo que durante cada franja horaria pueden ocurrir a lo sumo cuatro conciertos simultáneamente. 

    Se tiene un conjunto $B$ de bandas candidatas a ser invitadas a tocar en el festival. Para cada banda $i\in B$, se estima que atraerá $e_i$ espectadores a su show y se pronostica que dará una ganancia neta de $r_i$ pesos. Las bandas candidatas se categorizan en cinco conjuntos disjuntos según su género musical: $G_1, G_2, G_3, G_4, G_5$ (es decir: $\displaystyle\bigcup_{k=1}^5 G_k = B,\; G_k\cap G_r =\varnothing\;\text{si}\; k\neq r$).
    
    \begin{enumerate}
        \item Teniendo en cuenta que se desea decidir qué bandas invitar y en qué escenario y franja horaria tocará cada una de las bandas invitadas, modelar linealmente las siguientes restricciones del festival definiendo las variables y conjuntos que se consideren necesarios:
        
        \begin{enumerate}[a)]
            \item cada banda invitada debe tocar exactamente durante una franja horaria en exactamente un escenario.
            \item en cada franja horaria, en cada escenario puede tocar a lo sumo una banda.
            \item durante cada franja horaria pueden ocurrir a lo sumo cuatro conciertos simultáneamente.
            \item para cada franja horaria, la cantidad de espectadores del escenario $E_j$ no debe superar su capacidad $\ell_j$.
            \item la cantidad de espectadores totales en cada franja horaria no debe superar 20000.
            \item se deben invitar al menos $g_k$ bandas del género $G_k$.
            \item el escenario $4$ debe permanecer vacío durante la primera franja horaria de cada día.
            \item hay un conjunto $P\subset B$ de bandas muy populares que deben ser invitadas.
            \item cada banda del conjunto $P$ debe tocar durante alguna de las dos últimas franjas horarias de cualquiera de los días.
            \item la ganancia neta del festival debe superar los $R$ pesos.
            \item las bandas invitadas con menos de $S$ espectadores deben tocar durante los primeras tres franjas horarias.
        \end{enumerate}
        \item La empresa organizadora del festival está interesada en analizar el resultado obtenido con diferentes criterios. Utilizando las variables y conjuntos definidos en el ítem $1.$, modelar cada una de las siguientes funciones objetivo, agregando, de ser necesario, las restricciones y variables correspondientes.
        \begin{enumerate}[a)]
            \item Maximizar el mínimo de espectadores totales diarios.
            \item Minimizar el máximo número de espectadores totales durante la misma franja horaria.
            \item Minimizar la suma de la diferencia absoluta de cantidad de bandas invitadas entre cada par de géneros musicales.
        \end{enumerate}
    \end{enumerate}
\end{quest}

%%% TERMINA EL EJERCICIO 2 %%%

%%% EMPIEZA EL EJERCICIO 3 %%%


\begin{quest} 

    Una empresa forestal acaba de adquirir un territorio que desea explotar responsablemente. El terreno fue dividido en $N$ secciones y se desean producir $J$ tipos de madera distintos. Se cuenta con los siguientes datos:
    \begin{itemize}
        \item $m_j$: el árbol con madera $j$ tarda $m_j$ meses desde que fue plantado hasta su maduración (es decir, hasta estar listo para ser talado).
        \item $a_i$: área de la sección $i$.
        \item $A_t$: área máxima de tierra en la que se puede plantar árboles durante el mes $t$.
        % \item $vol_{ij\ell}$: volumen de madera $j$ obtenido al talar los árboles de la sección $i$ luego de $\ell$ meses de haber sido plantados.
        \item $v_{ij}$: volumen de madera $j$ que se obtiene al cosechar la sección $i$.
        % \item $d_{ij}$: cantidad de meses que debe dejarse descansar la tierra de la sección $i$ luego de talar los árboles con madera $j$.
        \item $g_j$: ganancia por cada unidad de volumen de madera $j$.
        \item $k_{ij}$: costo de plantar árboles de madera $j$ en la sección $i$. 
        \item $f_{it}$: costo de cada camión de transporte desde la sección $i$ hasta el aserradero en el mes $t$.
        \item $C_t$: cantidad de camiones de transporte disponibles durante el mes $t$.
        \item $V$: capacidad de volumen de madera de cada camión de transporte.
    \end{itemize}
    
    Cuando se cosecha la madera de una sección, se talan todos los árboles de la misma y todo el volumen de madera obtenido es transportado al aserradero. Cada camión de transporte hace a lo sumo un viaje por mes. El objetivo es planificar la producción de madera durante los próximos $\mc{T}$ meses. Plantear un modelo de programación lineal que permita decidir dónde, cuándo y qué árboles plantar y cuándo talarlos, de manera tal que se maximice la ganancia. Se deben considerar las siguientes restricciones:
    \begin{enumerate}
        \item en cada sección no puede haber más de una especie de árbol.
        \item en cada sección se puede plantar a lo sumo una vez.
        \item si se plantan árboles en una sección, deben ser talados en algún momento.
        \item no se puede talar árboles que no han madurado.
        \item en cada mes $t$, no se pueden plantar árboles en un área mayor que $A_t$.
        \item no sobrepasar la disponibilidad mensual de camiones
        \item no se puede plantar en las secciones del conjunto $\mc{S}$ durante los primeros $6$ meses.
        \item los árboles de la madera $4$ sólo pueden plantarse a partir del mes $12$, inclusive.
    \end{enumerate}
\end{quest}

%%% TERMINA EL EJERCICIO 3 %%%


%%% TERMINA EL PARCIAL %%%



\end{document}